% GENERATED FILE. Edit canonical lessons, then run npm run generate:books.
% canonical-source-hash: fnv1a64:fba3daef28b60861
% canonical-lessons: SA-C209-kavih, SA-C209-gayakah, SA-C209-nartakah, SA-C209-citrakarah, SA-C209-silpi

\chapter{Poet, Singer, Dancer, Painter, Craftsman}
\label{ch:sa-poet-singer-dancer-painter-craftsman}
\hlchaptermodality{209}

\begin{chapteropening}
\textbf{By the end of this chapter:} \emph{I can say a poet, a singer, a dancer, a painter and a craftsman.}

Complete the last lesson of chapter 209: I can say a poet, a singer, a dancer, a painter and a craftsman.
\end{chapteropening}

\section[\texorpdfstring{\sk{कविः} (kaviḥ)}{कविः (kaviḥ)}]{\sk{कविः} (kaviḥ) — a poet}

\label{lesson:SA-C209-kavih}



\begin{quote}
Before the new one: say the Sanskrit for an enemy, then the Sanskrit for a neighbour.
\end{quote}

\subsection*{You'll want to know: \sk{कविः}}
\textbf{\sk{कविः}} — \emph{kaviḥ} — \textquotedblleft{}a poet\textquotedblright{}.

\begin{grammarlens}[title={What you've built}]
One more word for this chapter.
\end{grammarlens}

\subsection*{Your turn}
\begin{itemize}
\raggedright
  \item \emph{Say it:} \emph{kaviḥ}
  \item \emph{Say it:} \emph{kaviḥ}, once more
  \item \emph{Say it:} say \emph{kaviḥ}
  \item \emph{Recall:} say the Sanskrit for an enemy, then the Sanskrit for a neighbour, then say \emph{kaviḥ} again
\end{itemize}

\begin{tcolorbox}[breakable,colback=teal!4,colframe=teal!35!black,title={Before you move on}]
What does \sk{कविः} mean? (\textquotedblleft{}A poet\textquotedblright{}.) Say it once more.
\end{tcolorbox}
\section[\texorpdfstring{\sk{गायकः} (gāyakaḥ)}{गायकः (gāyakaḥ)}]{\sk{गायकः} (gāyakaḥ) — a singer}

\label{lesson:SA-C209-gayakah}



\begin{quote}
Before the new one: say the Sanskrit for a neighbour, then the Sanskrit for a poet.
\end{quote}

\subsection*{You'll want to know: \sk{गायकः}}
\textbf{\sk{गायकः}} — \emph{gāyakaḥ} — \textquotedblleft{}a singer\textquotedblright{}.

\begin{grammarlens}[title={What you've built}]
One more word for this chapter.
\end{grammarlens}

\subsection*{Your turn}
\begin{itemize}
\raggedright
  \item \emph{Say it:} \emph{gāyakaḥ}
  \item \emph{Say it:} \emph{gāyakaḥ}, once more
  \item \emph{Say it:} say \emph{gāyakaḥ}
  \item \emph{Recall:} say the Sanskrit for a neighbour, then the Sanskrit for a poet, then say \emph{gāyakaḥ} again
\end{itemize}

\begin{tcolorbox}[breakable,colback=teal!4,colframe=teal!35!black,title={Before you move on}]
What does \sk{गायकः} mean? (\textquotedblleft{}A singer\textquotedblright{}.) Say it once more.
\end{tcolorbox}
\section[\texorpdfstring{\sk{नर्तकः} (nartakaḥ)}{नर्तकः (nartakaḥ)}]{\sk{नर्तकः} (nartakaḥ) — a dancer}

\label{lesson:SA-C209-nartakah}



\begin{quote}
Before the new one: say the Sanskrit for a poet, then the Sanskrit for a singer.
\end{quote}

\subsection*{You'll want to know: \sk{नर्तकः}}
\textbf{\sk{नर्तकः}} — \emph{nartakaḥ} — \textquotedblleft{}a dancer\textquotedblright{}.

\begin{grammarlens}[title={What you've built}]
One more word for this chapter.
\end{grammarlens}

\subsection*{Your turn}
\begin{itemize}
\raggedright
  \item \emph{Say it:} \emph{nartakaḥ}
  \item \emph{Say it:} \emph{nartakaḥ}, once more
  \item \emph{Say it:} say \emph{nartakaḥ}
  \item \emph{Recall:} say the Sanskrit for a poet, then the Sanskrit for a singer, then say \emph{nartakaḥ} again
\end{itemize}

\begin{tcolorbox}[breakable,colback=teal!4,colframe=teal!35!black,title={Before you move on}]
What does \sk{नर्तकः} mean? (\textquotedblleft{}A dancer\textquotedblright{}.) Say it once more.
\end{tcolorbox}
\section[\texorpdfstring{\sk{चित्रकारः} (citrakāraḥ)}{चित्रकारः (citrakāraḥ)}]{\sk{चित्रकारः} (citrakāraḥ) — a painter}

\label{lesson:SA-C209-citrakarah}



\begin{quote}
Before the new one: say the Sanskrit for a singer, then the Sanskrit for a dancer.
\end{quote}

\subsection*{You'll want to know: \sk{चित्रकारः}}
\textbf{\sk{चित्रकारः}} — \emph{citrakāraḥ} — \textquotedblleft{}a painter\textquotedblright{}.

\begin{grammarlens}[title={What you've built}]
One more word for this chapter.
\end{grammarlens}

\subsection*{Your turn}
\begin{itemize}
\raggedright
  \item \emph{Say it:} \emph{citrakāraḥ}
  \item \emph{Say it:} \emph{citrakāraḥ}, once more
  \item \emph{Say it:} say \emph{citrakāraḥ}
  \item \emph{Recall:} say the Sanskrit for a singer, then the Sanskrit for a dancer, then say \emph{citrakāraḥ} again
\end{itemize}

\begin{tcolorbox}[breakable,colback=teal!4,colframe=teal!35!black,title={Before you move on}]
What does \sk{चित्रकारः} mean? (\textquotedblleft{}A painter\textquotedblright{}.) Say it once more.
\end{tcolorbox}
\section[\texorpdfstring{\sk{शिल्पी} (śilpī)}{शिल्पी (śilpī)}]{\sk{शिल्पी} (śilpī) — a craftsman, a sculptor}

\label{lesson:SA-C209-silpi}



\begin{quote}
Before the new one: say the Sanskrit for a dancer, then the Sanskrit for a painter.
\end{quote}

\subsection*{You'll want to know: \sk{शिल्पी}}
\textbf{\sk{शिल्पी}} — \emph{śilpī} — \textquotedblleft{}a craftsman, a sculptor\textquotedblright{}.

\begin{grammarlens}[title={What you've built}]
One more word for this chapter.
\end{grammarlens}

\subsection*{Your turn}
\begin{itemize}
\raggedright
  \item \emph{Say it:} \emph{śilpī}
  \item \emph{Say it:} \emph{śilpī}, once more
  \item \emph{Say it:} say \emph{śilpī}
  \item \emph{Recall:} say the Sanskrit for a dancer, then the Sanskrit for a painter, then say \emph{śilpī} again
\end{itemize}

\begin{tcolorbox}[breakable,colback=teal!4,colframe=teal!35!black,title={Before you move on}]
What does \sk{शिल्पी} mean? (\textquotedblleft{}A craftsman, a sculptor\textquotedblright{}.) Say it once more.
\end{tcolorbox}
